\chapter{函数与模块}
\section{迭代器与生成器}
\subsection{迭代器}
\begin{lstlisting}[style=demostyle,numbers=none]
# 迭代器
# 迭代器是一个可以记住遍历的位置的对象
# 迭代器对象从集合的第一个元素开始访问，直到所有的元素被访问完结束
# 迭代器只能往前不会后退
# 迭代器有两个基本的方法：iter() 和 next()

list = [1, 2, 3, 4]
it = iter(list) # 创建迭代器对象
print(next(it)) # 输出迭代器的下一个元素
print(next(it))
\end{lstlisting}

\subsubsection{创建迭代器}
把一个类作为一个迭代器使用需要在类中实现两个方法 \codeinline{\_\_iter\_\_()} 与 \codeinline{\_\_next\_\_()}。
\begin{lstlisting}[style=pythonstyle]
class MyNumbers:
  def __iter__(self):
    self.a = 1
    return self
 
  def __next__(self):
    x = self.a
    self.a += 1
    return x
 
myclass = MyNumbers()
myiter = iter(myclass)
 
print(next(myiter))
print(next(myiter))
print(next(myiter))
print(next(myiter))
print(next(myiter))
\end{lstlisting}
\begin{lstlisting}[style=demostyle,numbers=none]
1
2
3
4
5    
\end{lstlisting}

\subsubsection{StopIteration异常}
StopIteration 异常用于标识迭代的完成，防止出现无限循环的情况，
在 \codeinline{\_\_next\_\_()} 方法中我们可以设置在完成指定循环次数后触发 StopIteration 异常来结束迭代。
\begin{lstlisting}[style=pythonstyle]
class MyNumbers:
  def __iter__(self):
    self.a = 1
    return self
 
  def __next__(self):
    if self.a <= 10:
      x = self.a
      self.a += 1
      return x
    else:
      raise StopIteration
 
myclass = MyNumbers()
myiter = iter(myclass)
 
for x in myiter:
  print(x)
\end{lstlisting}
\begin{lstlisting}[style=demostyle,numbers=none]
1
2
3
4
5
6
7
8
9
10   
\end{lstlisting}

\subsection{生成器}
生成器是一个返回迭代器的函数，只能用于迭代操作,
在调用生成器运行的过程中，每次遇到yield时函数会暂停并保存当前所有的运行信息，
返回yield的值，并在下一次执行\codeinline{next()}方法时从当前位置继续运行

\begin{lstlisting}[style=demostyle,numbers=none]
def fibonacci(n): # 生成器函数 - 斐波那契
    a, b, counter = 0, 1, 0
    while True:
        if counter > n:
            return
        yield a
        a, b = b, a + b
        counter += 1

f = fibonacci(10) # f 是一个迭代器，由生成器返回生成

for i in f:
    print(i, end=' ')
\end{lstlisting}


\section{Lambda表达式}
在很多场景，不太需要函数名字和返回值

\begin{lstlisting}[style=demostyle,numbers=none]
# Lambda表达式
# lambda [arg1 [,arg2,.....argn]]:expression

sum = lambda arg1, arg2: arg1 + arg2
print("相加后的值为：", sum(10, 20))
\end{lstlisting}


\begin{lstlisting}[style=demostyle,numbers=none]
G = 9.8  # 重力加速度常数
retrofire = lambda t, rtf, C: -C*G if t >= rtf else 0.0

# 等价于
def retrofire(t, rtf, C):
    if t >= rtf:
        return -C * G
    else:
        return 0.0
# 测试
print(retrofire(10, 5, 2)) # 输出 -19.6=-2*9.8
print(retrofire(3, 5, 2))  # 输出 0.0
\end{lstlisting}

\section{内建函数}
\begin{lstlisting}[style=demostyle,numbers=none]
# 内建函数
# abs() 返回数字的绝对值
print(abs(-10))

# map() 会根据提供的函数对指定序列做映射
def square(x):
    return x ** 2
print(list(map(square, [1, 2, 3, 4, 5])))
>>> [1, 4, 9, 16, 25]
a = [1, 2, 3, 4, 5]
print(list(map(lambda x: x+1 ,a)))
>>> [2, 3, 4, 5, 6]
b = [6, 7, 8, 9, 10]
print(list(map(lambda x, y: x + y, a, b)))
>>> [7, 9, 11, 13, 15]

# reduce() 函数会对参数序列中元素进行累积
## 在python3里面，reduce()函数需要导入functools模块
from functools import reduce 
def add(x, y):
    return x + y
print(reduce(add, [1, 2, 3, 4, 5]))
>>> 15

# filter() 函数用于过滤序列
## 在python3里面，filter()返回的是一个filter对象，需要用list()转换成列表
def is_odd(n):
    return n % 2 == 1
print(list(filter(is_odd, [1, 2, 3, 4, 5])))
>>> [1, 3, 5]


# sorted() 函数对所有可迭代的对象进行排序操作
print(sorted([1, 3, 2, 5, 4]))
>>> [1, 2, 3, 4, 5]
print(sorted([1, 3, 2, 5, 4], reverse=True))
>>> [5, 4, 3, 2, 1]

# zip() 函数用于将可迭代的对象作为参数，将对象中对应的元素打包成一个个元组
a = [1, 2, 3]
b = [4, 5, 6]
print(list(zip(a, b)))
>>> [(1, 4), (2, 5), (3, 6)]
## 字典值键互换
dict1 = {'a': 1, 'b': 2}
print(dict(zip(dict1.values(), dict1.keys())))
>>> {1: 'a', 2: 'b'}
\end{lstlisting}

\section{闭包}
函数嵌套函数，就是闭包,就是返回值是函数的函数

\begin{lstlisting}[style=demostyle,numbers=none]
# 不使用闭包
def addfunc():
    num1 = 2
    num2 = 3
    print(num1 + num2)

num = fun()

# 闭包
def addfunc1():
    num1 = 2
    def addfunc2(num2):
        print(num1 + num2)
    return addfunc2
    # addfunc2 表示一个函数的引用/函数名
    # addfunc2() 表示调用这个函数
num1 = addfunc1(3)  # 这里返回的是一个函数
num1(4)             # 这里调用的是返回的函数
>>> 6
\end{lstlisting}

\begin{lstlisting}[style=demostyle,numbers=none]
# 计算 a *x +b =y 的函数
# 计算 a * x + b = y
def func(a, b):
    def com_func(x):
        return a *x + b
    return com_func

fun = func(1, 4)   
print(fun(1))
print(fun(4))
print(fun(5))
\end{lstlisting}

\section{装饰器(Decorator)}
使用\codeinline{@}符号，把一个函数作为参数传递给另一个函数

闭包传进来的是一个变量，装饰器传进来的是一个函数
\begin{lstlisting}[style=demostyle,numbers=none]
# 装饰器
def decorator(func):
    def wrapper():
        print('调用函数')
        func()
    return wrapper

@decorator
def fun():
    print('hello world')

fun() # 相当于执行了 decorator(fun)()
\end{lstlisting}

\begin{lstlisting}[style=demostyle,numbers=none]
# 使用装饰器计算两个数的和
def tips(func):
    def nei(a, b):
        print('start')
        func(a, b)
        print('end')
    return nei

@tips  
def add(a, b):
    print(a + b)

add(1, 2) # 相当于执行了 tips(add)(1, 2)
\end{lstlisting}

\section{模块}
模块是代码量变得相当大时，为了将需要重复使用的有组织的代码放在一起，
方便维护和使用，这部分代码可以附加到其他代码中使用，
附加的过程就是\codeinline{import}。

\begin{lstlisting}[style=demostyle,numbers=none]
# 模块
import 模块名             # 导入整个模块
from 模块名 import 函数名 # 导入模块中的一个函数(不推荐，因为可能会覆盖掉原有的函数)
from 模块名 import *      # 导入模块中的所有函数
import 模块名 as 别名     # 导入模块并给模块取别名
\end{lstlisting}

\begin{lstlisting}[style=demostyle,numbers=none]
# 自定义模块
## @file mymodule.py
def add(a, b):
    return a + b

def sub(a, b):
    return a - b

# 同一文件夹下的test.py
## 如果不在同一文件夹下，需要将mymodule.py放到sys.path中的一个路径下
## @file test.py
import mymodule
print(mymodule.add(1, 2))
print(mymodule.sub(1, 2))
\end{lstlisting}




